% ============================ CHAPITRE 5 ==============================
\chapter{Segmentation des agences par apprentissage non supervisé}

\section*{Introduction}
Au-delà de la restitution descriptive offerte par le volet décisionnel, ce
chapitre introduit une brique d'\textbf{intelligence artificielle} : la
segmentation des agences par \textbf{apprentissage non supervisé} (\emph{clustering}).
L'objectif est de regrouper automatiquement les agences en profils homogènes,
sans étiquette préalable, afin d'éclairer les décisions (ciblage, allocation des
ressources, objectifs différenciés).

\section{Objectif et démarche}
La segmentation vise à \textbf{regrouper les agences} présentant des
caractéristiques d'activité similaires. La démarche suivie est la suivante :
préparation et \textbf{standardisation} des données, \textbf{détermination du
nombre de clusters} (méthode du coude et score de silhouette), \textbf{comparaison
de plusieurs modèles} de clustering, puis \textbf{interprétation} des segments
obtenus. L'implémentation s'appuie sur Python et la bibliothèque
\texttt{scikit-learn}.

\section{Données et variables}
Les données proviennent de l'agrégation, par agence, des vues décisionnelles
\texttt{V\_BI\_*}. Chaque agence est décrite par les variables du
tableau~\ref{tab:variables_clustering}.

\begin{table}[htbp]
\centering
\caption{Variables de segmentation (par agence)}
\label{tab:variables_clustering}
\renewcommand{\arraystretch}{1.2}
\begin{tabularx}{\textwidth}{>{\bfseries}p{4cm}X}
\toprule
\textmd{\textbf{Variable}} & \textbf{Description}\\
\midrule
effectif & Nombre d'employés de l'agence\\
nb\_gestionnaires & Nombre de gestionnaires\\
nb\_clients & Taille du portefeuille clients\\
total\_comptes & Total des ouvertures de comptes\\
production\_credits & Production de crédits\\
collecte\_epargne & Collecte d'épargne\\
taux\_presence & Taux de présence (issu du pointage)\\
\bottomrule
\end{tabularx}
\end{table}

\section{Prétraitement}
Les variables étant exprimées dans des unités et des ordres de grandeur très
différents (un effectif de quelques dizaines face à des milliers de comptes), une
\textbf{standardisation} (centrage-réduction, \texttt{StandardScaler}) est
appliquée. Sans cette étape, les variables de forte amplitude domineraient le
calcul des distances et fausseraient le regroupement.

\section{Détermination du nombre de clusters}
Le nombre de clusters $k$ n'étant pas connu a priori, il est déterminé en
combinant deux approches (figure~\ref{fig:choix_k}) : la \textbf{méthode du
coude}, qui repère le $k$ au-delà duquel l'inertie intra-cluster décroît peu, et
le \textbf{score de silhouette}, qui mesure la qualité de séparation des clusters.

\begin{figure}[htbp]
  \centering
  \includegraphics[width=\textwidth]{clustering_choix_k.png}
  \caption{Choix du nombre de clusters : méthode du coude et score de silhouette}
  \label{fig:choix_k}
\end{figure}

Les deux indicateurs concordent : l'inertie présente un \textbf{coude marqué à
$k=3$} et le score de silhouette y atteint son \textbf{maximum} ($\approx 0{,}71$).
La segmentation en \textbf{trois clusters} est donc retenue.

\section{Comparaison des modèles de clustering}
Quatre algorithmes ont été testés et comparés : \textbf{K-Means},
\textbf{classification ascendante hiérarchique} (agglomératif), \textbf{modèle de
mélange gaussien} (GMM) et \textbf{DBSCAN}. La comparaison s'appuie sur trois
indices internes : le \emph{score de silhouette} (plus il est élevé, mieux c'est),
l'indice de \emph{Davies-Bouldin} (plus il est faible, mieux c'est) et l'indice de
\emph{Calinski-Harabasz} (plus il est élevé, mieux c'est). Les résultats figurent
dans le tableau~\ref{tab:comparaison_clustering} et la
figure~\ref{fig:comparaison_clustering}.

\begin{table}[htbp]
\centering
\caption{Comparaison des modèles de clustering}
\label{tab:comparaison_clustering}
\renewcommand{\arraystretch}{1.25}
\begin{tabularx}{\textwidth}{>{\bfseries}lcccc}
\toprule
\textmd{\textbf{Modèle}} & \textbf{Nb clusters} & \textbf{Silhouette $\uparrow$} &
\textbf{Davies-Bouldin $\downarrow$} & \textbf{Calinski-Harabasz $\uparrow$}\\
\midrule
K-Means        & 3 & \textbf{0,709} & \textbf{0,402} & \textbf{276}\\
Agglomératif   & 3 & 0,709 & 0,402 & 276\\
GMM            & 3 & 0,709 & 0,402 & 276\\
DBSCAN         & 2 & 0,683 & 0,391 & 94\\
\bottomrule
\end{tabularx}
\end{table}

\begin{figure}[htbp]
  \centering
  \includegraphics[width=0.7\textwidth]{clustering_comparaison.png}
  \caption{Comparaison des modèles selon le score de silhouette}
  \label{fig:comparaison_clustering}
\end{figure}

K-Means, l'agglomératif et le GMM \textbf{convergent vers la même segmentation}
(silhouette de $0{,}709$), ce qui confirme la \textbf{robustesse} des trois
segments. DBSCAN, sensible à son paramètre de voisinage, ne distingue que deux
groupes. Le \textbf{K-Means} est retenu : il obtient le meilleur score, tout en
restant simple et facilement interprétable.

\section{Résultats et interprétation}
La figure~\ref{fig:clusters_pca} projette les agences dans le plan des deux
premières composantes principales (PCA), colorées selon leur cluster ; la première
composante (qui résume la \emph{taille} de l'agence) explique l'essentiel de la
variance. Les trois segments sont nettement séparés.

\begin{figure}[htbp]
  \centering
  \includegraphics[width=0.78\textwidth]{clustering_pca.png}
  \caption{Segmentation des agences (K-Means) projetée par PCA}
  \label{fig:clusters_pca}
\end{figure}

Le profil moyen de chaque segment (tableau~\ref{tab:profils_clusters}) permet de
les caractériser.

\begin{table}[htbp]
\centering
\caption{Profil moyen des segments d'agences}
\label{tab:profils_clusters}
\renewcommand{\arraystretch}{1.25}
\begin{tabularx}{\textwidth}{>{\bfseries}Xccccc}
\toprule
\textmd{\textbf{Segment}} & \textbf{Nb} & \textbf{Effectif} & \textbf{Clients} &
\textbf{Prod. crédits} & \textbf{Présence}\\
\midrule
Grandes agences   & 7  & 116 & 952 & 1752 & 0,90\\
Moyennes agences  & 16 & 55  & 427 & 789  & 0,90\\
Petites agences   & 22 & 22  & 184 & 337  & 0,80\\
\bottomrule
\end{tabularx}
\end{table}

On distingue ainsi : un segment de \textbf{grandes agences} (pôles majeurs à fort
effectif et forte production), un segment d'\textbf{agences moyennes} (activité
intermédiaire) et un segment de \textbf{petites agences} (effectif et activité
réduits, taux de présence plus faible).

\section{Apport décisionnel}
Cette segmentation constitue un véritable outil d'\textbf{aide à la décision} :
elle permet d'adapter les actions à chaque profil --- allouer les ressources selon
la taille, fixer des objectifs différenciés, accompagner spécifiquement les
petites agences ou capitaliser sur les grandes. Elle prolonge naturellement le
volet décisionnel et illustre l'apport de l'\textbf{intelligence artificielle} au
pilotage du réseau.

\section*{Conclusion}
Ce chapitre a montré comment l'apprentissage non supervisé enrichit la solution
d'une dimension analytique avancée. Après détermination du nombre optimal de
clusters et comparaison de plusieurs modèles, le K-Means a fait émerger trois
segments d'agences cohérents et exploitables pour la décision.
